problem:  
Let $\Sigma^{2n-1}$ be an exotic smooth sphere ($n \ge 4$).  
Fix the standard round metric $g_0$ on $S^{2n}$.  
Consider a sequence of smooth embeddings  
\[
\iota_j:\Sigma^{2n-1}\hookrightarrow (S^{2n},g_0)
\]
whose induced metrics $h_j = \iota_j^*g_0$ have nonnegative sectional curvature and whose second fundamental forms $A_j$ satisfy a uniform bound  
\[
\|A_j\|_{L^2} \le C
\]
for some fixed finite $C$, but **no pointwise curvature pinching is assumed**.  

Suppose moreover that $\iota_j(\Sigma^{2n-1})$ converges in the varifold or weak measure sense to an equatorial sphere $S^{2n-1}\subset S^{2n}$ as $j\to\infty$.  

**Problem:**  
Can such a sequence exist when $\Sigma^{2n-1}$ is exotic?  
Equivalently, can an exotic sphere appear as a curvature–controlled (in $L^2$ sense) limit of embedded hypersurfaces in the round sphere?  
If not, can one prove a quantitative **$L^2$‑umbilical rigidity theorem** stating that there exists $\varepsilon(n)>0$ so that whenever an embedded hypersurface in $(S^{2n},g_0)$ satisfies  
\[
\|A - \bar A\,\mathrm{Id}\|_{L^2} < \varepsilon(n)
\]
(for some mean curvature $\bar A$) and has nonnegative sectional curvature, it must be diffeomorphic to the standard sphere?

Why it is a “good’’ problem:  
This version eliminates near‑trivial rigidity by strengthening to a **weaker, integral extrinsic control**, placing the problem at the interface of geometric analysis and differential topology.  

* **Mathematical depth and simplicity:** The question asks whether weak extrinsic curvature bounds (in an averaged, global sense) still enforce smooth‑structure rigidity for hypersurfaces of the round sphere. It is simple to state but effectively open.  

* **Bridge between analytic geometry and topology:**  
  - From the **analytic side**, it relates to quantitative stability of the Simons identity, $L^2$‑umbilical inequalities, and varifold compactness (De Lellis–Müller type results).  
  - From the **topological side**, it connects to the classification of homotopy spheres and how far analytic curvature control can go toward determining smooth type. Resolving it would link smoothing theory with the fine analytic invariants of submanifold geometry.  

* **Conceptual impact:**  
  - A positive answer (prohibiting exotic limits under $L^2$ curvature control) would reveal a new **analytic smoothing rigidity**: nonnegative curvature plus weak extrinsic control would force the standard smooth structure.  
  - A negative answer (constructing exotic examples) would uncover a surprising **geometric flexibility** in weak curvature regimes, showing that exotic differential structures can persist even under analytic near‑umbilicity.

Its core strength is economy and foresight: a single, clear question—whether integral curvature control can detect exotic smooth structure—would, if answered, open a new path joining **global analysis, geometric measure theory, and differential topology**, fitting the archetype of a “good” advanced mathematical problem.